\subsection{A general three-state coned-bridge reduction}

\begin{lemma}[Proved: transitive pieces of a minimal-ideal core]
\label{CP2:A-act-minimal-core-orbits}
Let $M$ be a finite monoid, $K$ its minimal ideal, and $S$ a
nonempty finite right $M$-act. The core $U=SK$ is a disjoint
union of transitive right $K$-acts. Its singleton $K$-orbits
are exactly the states fixed by all of $M$. Every map
commuting with $M$ sends each core orbit onto a core orbit.
\end{lemma}
\begin{proof}
The finite minimal ideal is completely simple. In its Rees
description, $kK$ is a minimal right ideal for every $k\in K$,
and $k\ell K=kK$ for every $k,\ell\in K$. In particular
$k\ell t=k$ for some $t\in K$. If $u=sk$ and $v=u\ell$,
then $vt=u$. Also $u\in uK$, since $k\in kK$.
Thus $uK$ is transitive, and two such sets which intersect
are equal. These sets partition $SK$.

If $uK=\{u\}$, then for $m\in M$ and any $k\in K$,
$um=(uk)m=u(km)=u$. Conversely an $M$-fixed state belongs
to $SK$ and has singleton orbit. Finally, for an equivariant
map $f$, one has $f(uK)=f(u)K$, which is precisely the
core orbit containing $f(u)$.
\end{proof}

\begin{lemma}[Proved: a constant actor kills the first coned-bridge group]
\label{CP2:A-coned-connected-constant-any-center}
Let $S$ be a nonempty left-$N$, right-$M$ biset for finite
monoids $M,N$. If $Z=B(S,M)$ is connected and some left
$N$ actor on $S$ is constant, then
\[
 H_2\!\left(\operatorname{hocofib}
   (B(N,S,M)\longrightarrow BN);\mathbb Q\right)=0.
\]
No rational acyclicity assumption on $BM$ is required.
\end{lemma}
\begin{proof}
Let $C$ be the cone of the augmented left $\mathbb QN$
complex $C_*(Z;\mathbb Q)\to\mathbb Q$, with its ordinary
homological grading. Nonemptiness and connectedness give
$H_0(C)=H_1(C)=0$ and $H_2(C)=H_1(Z;\mathbb Q)$.
The derived $N$ coinvariants of this cone compute the
reduced chains of the displayed cofiber. The first possible
nonzero row of the hyperhomology spectral sequence gives
\[
 H_2(Y;\mathbb Q)=H_1(Z;\mathbb Q)_N.
\]
Indeed, in total degree two only its $(0,2)$ term occurs;
all possible incoming differentials have a row of degree
at most one and therefore vanish.

If an actor is constant with value $u$, commutation makes
$u$ an $M$-fixed state. Its map on $Z$ factors through the
forgetful map $Z\to BM$, followed by the fixed-state map
$BM\to Z$. Every finite monoid has finite universal group,
so $H_1(BM;\mathbb Q)=0$. This actor is consequently zero
on $H_1(Z;\mathbb Q)$. The coinvariant relation identifying
this zero image with each original class kills that group.
\end{proof}

\begin{theorem}[Proved: a three-state coned bridge reduces to an ordinary act]
\label{CP2:A-coned-three-transfer-to-ordinary}
Let $S$ be a nonempty left-$N$, right-$M$ biset with
$|S|\leq3$, for arbitrary finite monoids $M,N$, and put
\[
 Y=\operatorname{hocofib}(B(N,S,M)\longrightarrow BN).
\]
If $Y$ has the rational homology and cohomology ring of
$\CP^2$, then the right $M$ action on $S$ is trivial.
The ordinary left $N$-act cone on $S$ then has the same
rational homology and ring. If $Y\simeq\CP^2$, that
ordinary act cone is itself homotopy equivalent to $\CP^2$.
\end{theorem}
\begin{proof}
The binary case is excluded by
Lemma~\ref{CP2:A-star-coned-binary}, so assume $|S|=3$.
Use the core $U=SK$ of the preceding orbit lemma. The
commuting left action of $N$ preserves $U$ and sends core
orbits onto core orbits.

If $U$ has a unique core orbit, every $N$ actor restricts
to a surjection of this finite orbit onto itself, hence
to a permutation. The ordinary left action bar $B(N,U)$
maps to the two-sided bridge by using empty $M$ strings.
Its composite to $BN$ is the finite covering determined
by the permutation action on $U$. Rational transfer makes
$H^2(BN;\mathbb Q)\to H^2(B(N,U);\mathbb Q)$ injective,
so the same is true for pullback to the full bridge.
The cofiber exact sequence forces every second class of
$Y$ to vanish on $BN$. Such classes lift from the relative
group, whose quotient is a suspension. Their products are
zero, a contradiction.

If there are two core orbits and one has size two, the
other is the unique singleton orbit. Its state is fixed
by $M$, and equivariance forces $N$ to fix it as well.
Using that state and empty $M$ strings gives a section
$BN\to B(N,S,M)$. The same cohomology-injection argument
excludes this case.

All remaining core orbits are singletons, hence every
point of $U$ is fixed by all of $M$. If $U=S$, the right
action is trivial, as required. Otherwise the already
excluded unique-orbit case leaves exactly
\[
 U=\{u_0,u_1\},\qquad S=U\amalg\{w\}.
\]
Every right actor fixes both $u_i$, so its transformation
is the identity or one of the retractions $P_i$ sending
$w$ to $u_i$. At least one such retraction occurs, since
$U=SK\ne S$.

If only one $P_i$ occurs, it commutes with the entire
right action and, as an actual right actor, with every
left $N$ actor. It is a biset retraction onto two states.
Realize $Y$ by the source star with center $M$, leaf $N$
on $S$, and a trivial leaf on a singleton stalk. Leave
that singleton stalk unchanged and retract $S$ by $P_i$.
The rational retract-and-suspension argument of
Lemma~\ref{CP2:A-star-leaf-reductions}(4) transfers the
rational projective-plane target to this binary image,
contradicting the binary coned-bridge lemma.

If both $P_0,P_1$ occur, their common centralizer on $S$
is exactly the identity and the two full constant maps
to $u_0,u_1$. In fact, for a commuting map $f$,
$f(u_i)=P_i f(w)$; taking $f(w)=w,u_0,u_1$ gives exactly
those three possibilities. If every $N$ actor is the
identity, any state gives a section of the bridge
projection to $BN$, again killing the products. Otherwise
there is a constant $N$ actor. The right action space
$B(S,M)$ is connected, since the arrows representing
$P_0,P_1$ connect $w$ to both $u_i$. The preceding
constant-actor lemma makes $H_2(Y;\mathbb Q)=0$.
This finishes the proof that only the trivial right
action can occur.

In that final case replace the center $M$ by a trivial
corner, preserving both stalks and both leaves. This is
an actual category retract. Both connecting right actions
are trivial, so the opposite form of
Lemma~\ref{CP2:A-trivial-cut-H2} makes its pullback an
isomorphism on rational second cohomology. Its space is
the ordinary left $N$-act cone on $S$, with a contractible
whisker added at the cone vertex. As a retract, its homology
is bounded and is a direct summand of that of $Y$.
The second isomorphism and naturality of the nonzero square
therefore give exactly the rational projective-plane
homology and ring.

If $Y\simeq\CP^2$, use integral direct summands instead:
the retract has trivial fundamental group, and its only
possible positive homology groups are summands of
$\mathbb Z$ in degrees two and four. The rational second
isomorphism makes its integral second group $\mathbb Z$;
the split pullback to $H^2(Y;\mathbb Z)$ is then an
isomorphism. Naturality of the primitive square forces
the same conclusion in degree four, with a primitive
square on the retract. Integral projective-plane
recognition gives the asserted homotopy equivalence.
\end{proof}

\begin{corollary}[Proved: the leaf bound for a coned three-state bridge]
\label{CP2:A-coned-three-leaf-bounds}
For the preceding cofiber, $|N|\leq11$ excludes the
rational homology and cohomology ring of $\CP^2$.
The weaker hypothesis $|N|\leq17$ still excludes its
actual homotopy type. There is no restriction on $|M|$.
\end{corollary}
\begin{proof}
In the rational case the ordinary $N$-act cone supplied
by the theorem has bounded homology. Apply
Theorem~\ref{CP2:A-bounded-eleven-acts} and its opposite.
In the actual case the same ordinary cone has at most
$17+3+1=21$ arrows, contrary to
Theorem~\ref{CP2:A-three-state-minimum22}.
\end{proof}
